\uuid{crIE}
\exo7id{4726}
\titre{exo7 4726}
\auteur{quercia}
\organisation{exo7}
\datecreate{2010-03-16}
\isIndication{false}
\isCorrection{false}
\chapitre{Topologie}
\sousChapitre{Topologie de la droite réelle}
\cpge{mpsi-mp2i.nombres-reels-suites.a, pcsi.raisonnement-vocabulaire-ensembliste.c, ptsi.raisonnement-vocabulaire-ensembliste.c, mpsi-mp2i.nombres-reels-suites.c, pcsi.nombres-reels-suites.b, ptsi.nombres-reels-suites.b, mpsi-mp2i.nombres-reels-suites.d, pcsi.nombres-reels-suites.c, ptsi.nombres-reels-suites.c, bcpst1.suites-reelles}
\module{Analyse}
\niveau{L2}
\difficulte{}

\contenu{
\question{
Soit $(x_n)$ une suite r{\'e}elle born{\'e}e et $u>0,\,v>0$. On suppose que
$\frac{u}{v} \notin \Q$ et que les suites $(e^{iux_n})$ et $(e^{ivx_n})$
convergent.
Montrer que la suite $(x_n)$ converge.
}
}
